\section{Limitations and Future Work}
\label{sec:limitations}

Our study has several limitations. \textbf{Compute and ablations.} Because RLVR training is expensive, we only sweep the distillation learning-rate scale $\gamma$ and leave broader MinervaRL hyperparameter sweeps, including the distillation interval $I$, per-interval cap $M$, and EMA decay $\alpha$, for future work. \textbf{Data coverage.} Minerva-CTI and our evaluation suite are primarily English-centric; multilingual CTI, regional reporting styles, and non-English security artifacts remain important extensions. \textbf{Trace filtering.} MinervaRL uses lightweight heuristics and a TextCNN filter for ACR trace selection. Stronger LLM-based filters may improve trace quality but would increase cost, motivating more scalable quality estimators. \textbf{Training overhead.} MinervaRL adds overhead from ACR generation, log-probability computation, and periodic distillation. \hyperref[app:timing-overhead]{Appendix~\ref*{app:timing-overhead}} reports a 39.3\% increase in a 10-step Llama-3.1-8B timing study, although \Cref{fig:validation-accuracy-dynamics} shows faster progress toward strong validation performance. Reducing this overhead while preserving the support-seeding benefit is an important direction.